%% Đề 26 -- de-thi-thu-vao-lop-10-mon-toan-nam-2020-2021-phong-gddt-loc-binh-lang-son.pdf
\deheader{ĐỀ SỐ 26}%
  {Kỳ thi Tuyển sinh vào lớp 10 THPT — Năm học 2020--2021}%
  {Sở GD\&ĐT Lạng Sơn | Chỉ phần Đại số}
\addcontentsline{toc}{section}{Đề 26 — Tuyển sinh vào lớp 10 THPT — Năm học 2020--2021}

\begin{de_cau}{Câu 1}{2,5 điểm}
\begin{enumerate}[label=\alph*)]
\item Tính: $A = \sqrt{32}+3\sqrt{18}-\sqrt{50}$.
\item Rút gọn: $B=\dfrac{2\sqrt{x}-2}{\sqrt{x}-1}+\dfrac{\sqrt{x}-1}{\sqrt{x}+1}$, $x\geq 0, x\neq 1$.
\item Cho $C=\dfrac{x+2}{\sqrt{x}+1}-\sqrt{x}+1$. Tính $C$ khi $x=9$.
\end{enumerate}
\end{de_cau}

\begin{loi_giai_cau}
\textbf{a)} $4\sqrt{2}+9\sqrt{2}-5\sqrt{2}=8\sqrt{2}$.

\textbf{b)} $B=\dfrac{2(\sqrt{x}-1)}{\sqrt{x}-1}+\dfrac{\sqrt{x}-1}{\sqrt{x}+1}=2+\dfrac{\sqrt{x}-1}{\sqrt{x}+1}=\dfrac{2\sqrt{x}+2+\sqrt{x}-1}{\sqrt{x}+1}=\dfrac{3\sqrt{x}+1}{\sqrt{x}+1}$.

\textbf{c)} $x=9$: $C=\dfrac{11}{4}-3+1=\dfrac{11}{4}-2=\dfrac{3}{4}$.
\end{loi_giai_cau}

\begin{de_cau}{Câu 2}{2,5 điểm}
\begin{enumerate}[label=\alph*)]
\item Giải phương trình: $x^2+4x-5=0$.
\item Giải hệ phương trình: 
\[ \begin{cases}2x+3y=12\\x-2y=-2\end{cases} \].
\item Cho $(P)\colon y=x^2$ và $(d)\colon y=2x-1$. Tìm giao điểm của $(P)$ và $(d)$.
\end{enumerate}
\end{de_cau}

\begin{loi_giai_cau}
\textbf{a)} $(x+5)(x-1)=0$: $x=-5$ hoặc $x=1$.

\textbf{b)} Từ (2): $x=2y-2$. Thế (1): $2(2y-2)+3y=12\Rightarrow 7y=16\Rightarrow y=\frac{16}{7}, x=\frac{18}{7}$.

\textbf{c)} $x^2=2x-1\Rightarrow x^2-2x+1=0\Rightarrow(x-1)^2=0$: $x=1, y=1$. Tiếp xúc tại $(1;1)$.
\end{loi_giai_cau}

\begin{de_cau}{Câu 3}{1,5 điểm}
Hai người đi xe đạp từ A đến B. Người thứ nhất đi với vận tốc 15 km/h. Người thứ hai khởi hành sau
người thứ nhất 30 phút với vận tốc 20 km/h và đến B cùng lúc với người thứ nhất. Tính quãng đường AB.
\end{de_cau}

\begin{loi_giai_cau}
Gọi $AB=d$ (km). Thời gian người 1: $\dfrac{d}{15}$ giờ. Người 2 xuất phát sau $\frac{1}{2}$ giờ: thời gian $\dfrac{d}{20}$.

$\dfrac{d}{15} = \dfrac{d}{20}+\dfrac{1}{2}$.

$\dfrac{4d-3d}{60}=\dfrac{1}{2}\Rightarrow d=30$ km.
\end{loi_giai_cau}

\begin{de_cau}{Câu 4}{1,0 điểm}
Cho phương trình $x^2-2(m+1)x+m^2+2m=0$ (tham số $m$). Tìm $m$ để phương trình có hai nghiệm
$x_1, x_2$ thỏa mãn $x_1(x_1-2)+x_2(x_2-2)=2$.
\end{de_cau}

\begin{loi_giai_cau}
$\Delta'=(m+1)^2-(m^2+2m)=1>0$ với mọi $m$. Phương trình luôn có 2 nghiệm phân biệt.

Vi-ét: $S=2(m+1)$, $P=m^2+2m=m(m+2)$.

Theo giả thiết:
\[x_1(x_1-2)+x_2(x_2-2) = x_1^2+x_2^2-2(x_1+x_2) = S^2-2P-2S.\]
Thay $S = 2(m+1)$ và $P = m(m+2)$:
\[S^2-2P-2S = 4(m+1)^2-2m(m+2)-4(m+1) = 2m^2.\]
Theo đề bài: $2m^2 = 2 \Leftrightarrow m^2 = 1 \Leftrightarrow \boxed{m=\pm 1}$.
\end{loi_giai_cau}
